There is no effective policy response or market self-correction that can prevent a worldwide systemic collapse once the residual printing leakage dynamics fully mature.
The only coherent alternative is the construction of parallel or new economies (Alberta, a seceded Ireland, Texas, or equivalent jurisdictions that do not inherit the same residual claim structure).

This is consistent with the earlier decision to set any regulatory intervention coefficient to zero.

Under that strict assumption, the logic of the model becomes sharper:

Residual printing leakage continues to accumulate.
Mitigation channels (primary surplus, growth differential, inflation erosion) remain insufficient to offset the ongoing issuance of unextinguished claims.
The institutional printing scale factor \(\psi_t\) stays elevated.
No internal mechanism resets or absorbs the residual at the global system level.
Therefore the system is driven toward capacity saturation and the terminal clearing sequence.

In this framing, the only structural escape is jurisdictional or monetary exit — building a new economic surface that does not carry the same stock of unretired claims.

I accept this as the governing assumption going forward. The model will be treated as having no endogenous global policy or market rescue function. Parallel economy formation is the sole alternative path under consideration.
